\documentclass{beamer}

% 
	%! Author = Matthew
%! Date = 31.03.2023

% Preamble
\usepackage{xcolor}
\usepackage{hyperref}
\usepackage{comment}
\usepackage[]{ocgx2}
\usepackage[utf8]{inputenc}
\usepackage[russian]{babel}
\usepackage[T1]{fontenc}
\usepackage{xfrac}

\usepackage[colorinlistoftodos]{todonotes}
\usepackage{amsfonts}
\usepackage{amsmath}
\usepackage{amssymb}
\usepackage{amsthm}
\usepackage{tikz} 
\usepackage{indentfirst}
\usepackage{mathtools}
\usepackage{nccmath}
\usepackage{transparent}
\usepackage{mathtools}
\usepackage{graphics}
\usepackage{physics}


\hypersetup{%
  colorlinks=false,
  linkbordercolor=cyan
}

\usepackage{upgreek}

%вставка изображений и картинок всяких

\usepackage{import}
\usepackage{xifthen}
\usepackage{pdfpages}
\usepackage{transparent}

\newcommand{\incfig}[1]{%
    \def\svgwidth{\columnwidth}
    \import{./pictures/}{#1.pdf_tex}
}
\graphicspath{{pictures/}}
\DeclareGraphicsExtensions{.pdf,.png,.jpg, .jpeg, .tex, .pdf_tex}

\usepackage{booktabs} % для таблиц
\usepackage{enumitem} % для списков


\usepackage{tikz-cd} % коммутативные диаграммы

% Шрифты
\usepackage[english]{babel}
\usepackage[T1]{fontenc}

\usepackage{titling}
\usepackage{textcomp}% специальные символы в тексте

\usepackage{mathtext}
\usepackage{amsmath,amsfonts,amssymb,amsthm,mathtools} % математика
\usepackage{icomma} % умная запятая
\usepackage{import} %  импортирование


\usepackage{pdfpages} % мультри-пдф страницы
\usepackage{transparent} % что-то про цвета
\usepackage{xcolor} % цвета

\usepackage{caption} % комментарии к figure
\usepackage{epigraph} % эпиграфы

\usepackage{comment} % удобные комментарии
\usepackage{xfrac} % дроби
\usepackage{moresize} % все размеры шрифтов
\usepackage{dsfont} % mathbb для всего

\usepackage{yfonts} % готические буквы

% Окружения для математики:

\newtheorem{statement}{Предложение}
\newtheorem{corollary}{Следствие}
\newtheorem{theorem}{Теорема}
\newtheorem{hypothesis}{Гипотеза}
\theoremstyle{definition}
\newtheorem{definition}{Определение}
\newtheorem*{*definition}{Определение}
\newtheorem{example}{Пример}
\newtheorem{antiexample}{Антипример}
\newtheorem{lemma}{Лемма}
\theoremstyle{remark}
\newtheorem{remark}{Замечание}
\newtheorem*{*remark}{Замечание}
\newtheorem*{exercise}{Упражнение}

% Настройки счетчиков:

%\numberwithin{equation}{section} % Number equations within sections (i.e. 1.1, 1.2, 2.1, 2.2 instead of 1, 2, 3, 4)
%\numberwithin{figure}{section} % Number figures within sections (i.e. 1.1, 1.2, 2.1, 2.2 instead of 1, 2, 3, 4)
%\numberwithin{table}{section} % Number tables within sections (i.e. 1.1, 1.2, 2.1, 2.2 instead of 1, 2, 3, 4)

% Геометрия файла:


\usepackage[left=2.5cm,right=2.5cm,top=2.5cm,bottom=3.5cm]{geometry}

\setlist{noitemsep} % No spacing between list items

%\geometry{left=1.5cm,right=1.5cm,top=2.5cm,bottom=2cm, a4paper}
% Счётчики разделов:


%\sectionfont{\vspace{6pt}\centering\normalfont\scshape} % \section{} styling
%\subsectionfont{\normalfont\bfseries} % \subsection{} styling
%\subsubsectionfont{\normalfont\itshape} % \subsubsection{} styling
%\paragraphfont{\normalfont\scshape} % \paragraph{} styling

%\newcommand{\RNumb}[1]{\uppercase\expandafter{\romannumeral #1\relax}}

%\renewcommand\thesection{\arabic{section}.}
%\renewcommand\thesubsection{\thesection\arabic{subsection}}
%\renewcommand\thesubsubsection{\RNumb{\arabic{subsubsection}}}
%\renewcommand{\bf}{\textbf}

% Колонтитулы

%\usepackage{fancyhdr}
%\pagestyle{fancy}
%\fancyhf{} % clear all fields
%\fancyhead[L]{\rightmark}
%\fancyhead[R]{\thepage}

%\renewcommand{\sectionmark}[1]{%
%  \markright{\thesection\ #1}}%
%\setlength{\headheight}{17.0pt}
%\addtolength{\topmargin}{-2.49998pt}



% Операторы:
\newcommand{\id}{\operatorname{id}}
% \newcommand{\rank}{\operatorname{rank}}
\DeclareMathOperator{\ord}{ord}
\DeclareMathOperator{\ld}{ld}
\DeclareMathOperator{\exi}{exi}
\DeclareMathOperator{\num}{num}
\DeclareMathOperator{\card}{Card}
\DeclareMathOperator{\sk}{sk}
\DeclareMathOperator{\den}{den}
\DeclareMathOperator{\ran}{ran}
\DeclareMathOperator{\dom}{dom}
\DeclareMathOperator{\diam}{diam}
\DeclareMathOperator{\Aut}{Aut}
\DeclareMathOperator{\Par}{P}
\DeclareMathOperator{\St}{St}
\DeclareMathOperator{\Div}{div}
\DeclareMathOperator{\sign}{sign}
\DeclareMathOperator{\Int}{Int}
\DeclareMathOperator{\RelInt}{RelInt}
\DeclareMathOperator{\Val}{Val}
\DeclareMathOperator{\WP}{WP}
\DeclareMathOperator{\Det}{det}
\DeclareMathOperator{\argmin}{argmin}

\DeclareMathOperator{\Cl}{Cl}
\DeclareMathOperator{\pr}{pr}
\DeclareMathOperator{\supp}{supp}
\DeclareMathOperator{\Af}{Aff}
\renewcommand{\Im}{\mathop{\mathrm{Im}}\nolimits}
\DeclareMathOperator{\conv}{conv}
\DeclareMathOperator{\Fr}{Fr}
\DeclareMathOperator{\Span}{span}
\DeclareMathOperator{\Map}{Map}
\DeclareMathOperator{\Ker}{Ker}
\DeclarePairedDelimiter\lr{(}{)}
\DeclareRobustCommand{\divby}{%
     \mathrel{\text{\vbox{\baselineskip.65ex\lineskiplimit0pt\hbox{.}\hbox{.}\hbox{.}}}}%
}
\newcommand{\eqdef}{\stackrel{\mathrm{def}}{=}}
\DeclareRobustCommand{\notdivby}{%
     \!\!\not\;\divby%
}
\newcommand{\Mod}[1]{\ (\mathrm{mod}\ #1)}


% Tikz и графика:


\usepackage{pgfplots}
\usepackage{tikz}
\usetikzlibrary{3d,perspective}
\usetikzlibrary{animations}
%\usepackage{tikz-cd} % коммутативные диаграммы
\usetikzlibrary{cd, shapes.geometric}
\pgfplotsset{width=6cm,compat=newest}


%%%% гиперссылки

%\usepackage[colorlinks, urlcolor = blue, filecolor = blue,citecolor = blue, linkcolor = blue]{hyperref}

% Буковы
\def\<{\langle}
\def\>{\rangle}
\newcommand{\Lin}{\operatorname{Lin}}
\newcommand{\N}{\mathbb{N}}
\newcommand{\Z}{\mathbb{Z}}
\newcommand{\Q}{\mathbb{Q}}
\newcommand{\R}{\mathbb{R}}
%%\let\oldC\C
%\newcommand{\C}{\mathbb{C}}
\newcommand{\F}{\mathbb{F}}

\let\oldAA\AA
\renewcommand{\AA}{\mathbb{A}}
\newcommand{\DD}{\mathbb{D}}
\newcommand{\ga}{\gamma}
\newcommand{\co}{\colon}
\newcommand{\ep}{\varepsilon}
\newcommand{\la}{\lambda}
\newcommand{\ka}{\kappa}
\newcommand{\EE}{\mathbb{E}}
\newcommand{\HH}{\mathbb{H}}
\newcommand{\KK}{\mathbb{K}}
\newcommand{\OO}{\mathbb{O}}
\newcommand{\PP}{\mathbb{P}}
\let\oldSS\SS
\renewcommand{\SS}{\mathbb{S}}
\newcommand{\TT}{\mathbb{T}}

\newcommand{\cA}{\mathcal{A}}
\newcommand{\cB}{\mathcal{B}}
\newcommand{\cC}{\mathcal{C}}
\newcommand{\cD}{\mathcal{D}}
\newcommand{\cE}{\mathcal{E}}
\newcommand{\cF}{\mathcal{F}}
\newcommand{\cG}{\mathcal{G}}
\newcommand{\cH}{\mathcal{H}}
\newcommand{\cI}{\mathcal{I}}
\newcommand{\cJ}{\mathcal{J}}
\newcommand{\cK}{\mathcal{K}}
\newcommand{\cL}{\mathcal{L}}
\newcommand{\cM}{\mathcal{M}}
\newcommand{\cN}{\mathcal{N}}
\newcommand{\cO}{\mathcal{O}}
\newcommand{\cQ}{\mathcal{Q}}
\newcommand{\cP}{\mathcal{P}}
\newcommand{\cR}{\mathcal{R}}
\newcommand{\cS}{\mathcal{S}}
\newcommand{\cT}{\mathcal{T}}
\newcommand{\cU}{\mathcal{U}}
\newcommand{\cV}{\mathcal{V}}
\newcommand{\cW}{\mathcal{W}}
\newcommand{\cX}{\mathcal{X}}
\newcommand{\cY}{\mathcal{Y}}
\newcommand{\cZ}{\mathcal{Z}}

\newcommand{\rD}{\mathrm{D}}
\newcommand{\rK}{\mathrm{K}}
\newcommand{\rP}{\mathrm{P}}
\newcommand{\rT}{\mathrm{T}}

\newcommand{\fA}{\mathfrak{A}}
\newcommand{\fB}{\mathfrak{B}}
\newcommand{\fM}{\mathfrak{M}}
\newcommand{\fL}{\mathfrak{L}}
\newcommand{\fX}{\mathfrak{X}}

\newcommand{\fm}{\mathfrak{m}}
\newcommand{\fp}{\mathfrak{p}}

\newcommand{\RomanNumeralCaps}[1]
{\MakeUppercase{\romannumeral #1}}
\newcommand{\II}{\textrm{\RomanNumeralCaps{2}}}
\newcommand{\I}{\textrm{\RomanNumeralCaps{1}}}

\usepackage{csquotes}

\usetheme{Madrid}
\usepackage[utf8]{inputenc}
\usepackage[russian]{babel}
\usepackage{amsmath, amssymb, bm}
\usepackage{animate}
\usepackage{graphicx}

\title{Flow Matching. Quick Tour}
\author{}
\date{}

\DeclareMathOperator{\argmin}{argmin}

\begin{document}

\frame{\titlepage}

%------------------------------------------------
\section{Основные идеи}

%------------------------------------------------
\begin{frame}{Постановка задачи генерации}

Пусть:
\begin{itemize}
    \item $p_0(x)$ — простое распределение (например, $\mathcal N(0,I)$),
    \item $p_1(x)$ — целевое распределение данных.
\end{itemize}
Задача: построить диффеоморфизм
\[
\psi_1 : \mathbb R^d \to \mathbb R^d,
\]
такой что
\[
p_1(x) = [\psi_1]_\# p_0 (x) = p_0(\psi_1^{-1}(x)) |\partial_x \psi_1^{-1}(x)|
\]

Т.е. $X_1 = \psi_1(X_0) \sim p_1 \iff X_0 \sim p_0$ 

\end{frame}

\begin{frame}{Постановка задачи генерации}

Вместо прямого построения $\psi_1$ рассматриваем поток
\[
\frac{d}{dt}\psi_t(x) = v_t(\psi_t(x)), \quad t \in [0,1].
\]
Из которого отображение $\psi_1$ восстанавливается интегрированием
\[
\psi_1(x) = p_0(x) + \int_0^{1} v_t(\psi_t(x)) dt
\]
\end{frame}

%------------------------------------------------
\begin{frame}{Поток и уравнение непрерывности}

Если $\psi_t$ удовлетворяет
\[
\frac{d}{dt}\psi_t(x) = v_t(\psi_t(x)),
\]
то плотность $p_t = [\psi_t]_\# p_0 (x)$ удовлетворяет уравнению непрерывности:
\[
\partial_t p_t(x) + \nabla \cdot \big(p_t(x) v_t(x)\big) = 0.
\]

Тогда, если мы найдём правильное $v_t$, то $\psi_t$ \textit{хорошим образом} перенесёт $p_0$ в $p_1$

\end{frame}

\begin{frame}
    \begin{figure}
        \centering
        \animategraphics[loop, autoplay, width=0.8\linewidth]{30}{anim/FM/frame_}{0001}{0150}
    \end{figure}
\end{frame}



%------------------------------------------------
\begin{frame}{Ключевая идея FM}

Вместо обучения напрямую распределения $q_\theta(x) \approx p_1(x)$ через KL-дивергенцию,
Flow Matching предлагает:

\medskip
\begin{enumerate}
    \item Задать интерполирующее распределение $p_t$ между $p_0$ и $p_1$
    \item Вычислить соответствующее \textit{истинное} поле скоростей $u_t(x)$.
    \item Обучать модель $v_\theta$ по регрессии:
\end{enumerate}

\[
\mathcal L_{\text{FM}} =
\mathbb E_{t \sim U[0,1],\, x \sim p_t}
\| v_t^\theta(x) - u_t(x) \|^2.
\]

\end{frame}

\begin{frame}
    
    \centering
    \begin{tabular}{c c}
        \includegraphics[width=0.35\linewidth]{images/data.jpeg} & \includegraphics[width=0.35\linewidth]{images/path_design.jpeg} \\
        \includegraphics[width=0.35\linewidth]{images/training.jpeg} & \includegraphics[width=0.35\linewidth]{images/sampling.jpeg} \\
    \end{tabular}
    
    Проблема лишь в том, что \textit{истинное} поле $u_t(x)$ нам обычно неизвестно
\end{frame}



%------------------------------------------------
\section{От conditional к marginal loss}

%------------------------------------------------
\begin{frame}{Условная формулировка}

Вместо работы с маргинальным $u_t(x)$ вводится

\[
u_t(x \mid z),
\]

где $z$ — латентная случайная величина (например, пара $z = (x_0, x_1) \sim \pi_{0, 1} $).

Тогда условный лосс:
\[
\mathcal L_{\text{cond}} =
\mathbb E_{t,z,x_t \sim p_{1|z}(\cdot \mid z)}
\| v_t^\theta(x_t) - u_t(x_t \mid z) \|^2.
\]

Это легко сэмплировать:
\begin{enumerate}
    \item берём $z$,
    \item генерируем $x_t \sim p_{1 | z}$,
    \item знаем $u_t(x_t \mid z)$ аналитически.
\end{enumerate}

Тогда истинное поле выражается как

\[
u_t(x)
=
\mathbb{E}_{z \sim \pi} \big[
u_t(X_t \mid z) \mid X_t = x
\big].
\]

\end{frame}

%------------------------------------------------
\begin{frame}{Пример: линейная интерполяция}

Простейший случай:

\begin{itemize}
    \item $p_0(x) = \mathcal N(\mu_0, \sigma_0)$
    \item $p_1(x) = \mathcal N(\mu_1, \sigma_1)$
    \item $\psi_{t \mid 1}(x \mid x_1) = t x_1 + (1 - t) x$
\end{itemize}

Тогда условное поле скоростей:
\[
u_{t \mid 1}(x \mid x_1) = \dot{\psi}_{t \mid 1} \circ \psi_{t \mid 1}^{-1}(x \mid x_1) = \frac{x_1 - x}{1 - t}
\]

И условный probability path:

\[
p_{t \mid 1} (x \mid x_1) = p_{0 \mid 1} (\psi_{t \mid 1}^{-1}(x \mid x_1)) | \partial_x\psi_{t \mid 1}^{-1}(x \mid x_1) | = 
\]
\[
= \frac{1}{\sqrt{2 \pi} \sigma} e^{-\frac{1}{2} \Big( \frac{\frac{x - t x_1}{1 - t} - \mu_0}{\sigma_0} \Big)^2} \cdot \frac{1}{1 - t} = 
\mathcal N(t x_1 + (1 - t)\mu_0, (1 - t)\sigma_0)
\]

\end{frame}

\begin{frame}{Пример: линейная интерполяция}
    Тогда marginal probability path:
    \[
    p_t(x) = \int p_{t \mid 1}(x | x_1) p_1(x_1) dx_1 = 
    \]
    \[
    = \int \mathcal N(t x_1 + (1 - t)\mu_0, (1 - t)\sigma_0) \> \mathcal N(\mu_1, \sigma_1) dx_1 =
    \]
    \[
    = \mathcal{N}(t \mu_1 + (1 - t) \mu_0, t^2 \sigma_1 + (1 - t) \sigma_0)
    \]
    Отсюда по формуле Байеса 
    \[
    p_{1 \mid t}(x_1 \mid x) = \frac{p_{t \mid 1}(x \mid x_1) p_1(x_1)}{p_t(x)} = \mathcal N(\mu_{1 | t}(x), \sigma_{1|t}(x))
    \]
    где
    \[
    \mu_{1 \mid t}(x) = \mu_1 + \frac{t \sigma_1 (x - t\mu_1 - (1 - t)\mu_0)}{t^2 \sigma_1 + (1 - t)\sigma_0}
    \]
    \[
    \sigma_{1 \mid t}(x) = \frac{\sigma_1 (1 - t) \sigma_0}{t^2 \sigma_1 + (1 - t)\sigma_0}
    \]
\end{frame}

\begin{frame}
    Итого, marginal velocity field:
    \[
    u_t(x) = \int u_{t \mid 1}(x \mid x_1) p_{1 \mid t}(x_1 \mid x) dx_1 = 
    \]
    \[
    = \mathbb{E}\Big[\frac{X_1 - X_t}{1 - t} \mid X_t = x\Big] = \frac{1}{1 - t}\mathbb{E}[X_1 \mid X_t = x] - \frac{x}{1 - t} =
    \]
    \[
    = \frac{\mu_{1 | t}(x) - x}{1 - t}
    \]
\end{frame}

\begin{frame}
    
    \centering
    \begin{tabular}{c c}
        \includegraphics[width=0.3\linewidth]{images/cond_path.jpeg} & \includegraphics[width=0.3\linewidth]{images/marg_path.jpeg} \\
        \small{$\mathcal{N}(t x_1 + (1 - t)\mu_0, (1 - t)\sigma_0)$} & \small{$\mathcal{N} (t \mu_1 + (1 - t) \mu_0, t^2 \sigma_1 + (1 - t) \sigma_0)$} \\
        \includegraphics[width=0.3\linewidth]{images/cond_vel.jpeg} & \includegraphics[width=0.3\linewidth]{images/marg_vel.jpeg} \\
        \Large{$\frac{x_1 - x}{1 - t}$} & \Large{$\frac{\mu_{1 | t}(x) - x}{1 - t}$} \\
    \end{tabular}

\end{frame}

% \begin{frame}
%     \begin{figure}
%         \centering
%         \animategraphics[loop, autoplay, width=0.8\linewidth]{30}{anim/CFM/frame_}{0001}{0150}
%     \end{figure}
% \end{frame}


%------------------------------------------------
\begin{frame}{Теорема о совпадении лоссов}

\begin{block}{Теорема}

Градиент условного лосса совпадает
с градиентом маргинального лосса:

\[
\nabla_\theta \mathcal L_{\text{cond}}
=
\nabla_\theta \mathcal L_{\text{marg}},
\]

где

\[
\mathcal L_{\text{cond}} =
\mathbb E_{t,Z \sim \pi, X_t \sim p_{1 | Z}}
\| v_t^\theta(X_t) - u_t(X_t|Z) \|^2.
\]

\[
\mathcal L_{\text{marg}}
=
\mathbb E_{t, X_t \sim p_t}
\| v_t^\theta(X_t) - u_t(X_t) \|^2
\]

\end{block}

\end{frame}

\begin{frame}{Доказательство}
    \begin{align*}
        \nabla_\theta \mathcal L_{\text{marg}} 
        & = \nabla_\theta \mathbb{E}_{t,X_t} \| v_t^\theta(X_t) - u_t(X_t) \|^2 \\
        & = \mathbb{E}_{t,X_t} \nabla_\theta \| v_t^\theta(X_t) - u_t(X_t) \|^2 \\
        & = \mathbb{E}_{t,X_t} \textcolor{blue}{\nabla_\theta v_t^\theta(X_t) \nabla_{v_t^\theta}} \| v_t^\theta(X_t) - u_t(X_t) \|^2 \\
        & = \mathbb{E}_{t,X_t} \nabla_\theta v_t^\theta(X_t) \nabla_{v_t^\theta} \| v_t^\theta(X_t) - \textcolor{blue}{\mathbb{E}_{Z \sim p_{Z | t}} [ u_t(X_t \mid Z)]} \|^2 \\
        & = \mathbb{E}_{t,X_t} \textcolor{blue}{\mathbb{E}_{Z \sim p_{Z | t}}} \nabla_\theta v_t^\theta(X_t) \nabla_{v_t^\theta} \| v_t^\theta(X_t) - u_t(X_t \mid Z) \|^2 \\
        & = \mathbb{E}_{t,X_t} \textcolor{blue}{\mathbb{E}_{Z \sim p_{Z | t}}} \nabla_\theta \| v_t^\theta(X_t) - u_t(X_t \mid Z) \|^2 \\
        & = \mathbb{E}_{t, \textcolor{blue}{X_t \sim p_{t|Z}, Z \sim \pi}} \nabla_\theta \| v_t^\theta(X_t) - u_t(X_t \mid Z) \|^2 \\
        & = \nabla_\theta \mathbb{E}_{t, X_t \sim p_{t|Z}, Z \sim \pi} \| v_t^\theta(X_t) - u_t(X_t \mid Z) \|^2 \\
        & = \nabla_\theta \mathcal{L}_{\text{cond}} 
    \end{align*}

\end{frame}



%------------------------------------------------
% \begin{frame}{Главная идея доказательства}

% Рассмотрим разложение:

% \[
% u_t(x|z)
% =
% u_t(x)
% +
% \varepsilon,
% \]

% где

% \[
% \mathbb E[\varepsilon \mid x] = 0.
% \]

% Тогда:

% \[
% \|v - u_t(x|z)\|^2
% =
% \|v - u_t(x)\|^2
% +
% \|\varepsilon\|^2
% +
% 2\langle v-u_t(x), \varepsilon \rangle.
% \]

% После взятия условного ожидания:

% \[
% \mathbb E[\langle v-u_t(x), \varepsilon \rangle \mid x] = 0.
% \]

% \end{frame}

%------------------------------------------------
\begin{frame}{Следствие}

Получаем:

\[
\mathcal L_{\text{cond}}
=
\mathcal L_{\text{marg}}
+
\text{const}.
\]

Т.е. их $\argmin$ совпадают, итого: 

\begin{block}{Следствие}
    Зная только \textbf{условное} поле скоростей, мы можем выучить \textbf{истинное}
\end{block}

\end{frame}

\begin{frame}

    \centering
    \includegraphics[width=\linewidth]{images/fm_scheme.jpeg}

\end{frame}


%------------------------------------------------
\begin{frame}{Итог}

Flow Matching:

\begin{itemize}
    \item хорошая лосс-функция (в отличие от KL-дивергенции)
    \item позволяет обучать настоящий поток через условный
    \item ODE солвер используется только на этапе семплирования, но не обучения
\end{itemize}
\end{frame}

%------------------------------------------------
\end{document}